\section{Phụ lục}

\subsection{Chứng minh Bổ đề tập trung phân phối $\chi^2$}
\label{app:chi2-proof}

Bổ đề~\ref{lem:chi2-conc} được dùng trong chứng minh bổ đề Johnson--Lindenstrauss (Mục~\ref{sec:jl}). Dưới đây là chứng minh đầy đủ bằng kỹ thuật Chernoff bound.

\begin{proof}[Chứng minh Bổ đề~\ref{lem:chi2-conc}]
Dùng bất đẳng thức Chernoff (Markov trên $e^{tS}$). Với $t < \frac{1}{2}$, MGF của $Z_i^2$ là $\E[e^{tZ_i^2}] = (1-2t)^{-\frac{1}{2}}$, nên $\E[e^{tS}] = (1-2t)^{-\frac{k}{2}}$.

\textbf{Đuôi trên} ($S > (1+\varepsilon)k$): với $t \in \left(0, \frac{1}{2}\right)$,
\[
\Pr(S > (1+\varepsilon)k) \leq \frac{(1-2t)^{-k/2}}{e^{t(1+\varepsilon)k}}.
\]
Chọn $t = \frac{\varepsilon}{2(1+\varepsilon)}$ và dùng $\ln(1+\varepsilon) \leq \varepsilon - \frac{\varepsilon^2}{4}$ (với $\varepsilon \in (0,1)$) suy ra $\Pr(S > (1+\varepsilon)k) \leq e^{-\frac{\varepsilon^2 k}{8}} \leq e^{-\frac{\varepsilon^2 k}{4}}$.

\textbf{Đuôi dưới} ($S < (1-\varepsilon)k$): với $t \in \left(-\frac{1}{2}, 0\right)$, chọn $t = -\frac{\varepsilon}{2(1-\varepsilon)}$ và $\ln(1-\varepsilon) \leq -\varepsilon - \frac{\varepsilon^2}{2}$ cho $\Pr(S < (1-\varepsilon)k) \leq e^{-\frac{\varepsilon^2 k}{4}}$.

Union bound trên hai đuôi: $\Pr\!\left(\left|\frac{S}{k} - 1\right| > \varepsilon\right) \leq 2 e^{-\frac{\varepsilon^2 k}{4}}$.
\end{proof}

\subsection{Khai triển gradient UMAP}
\label{app:umap-gradient}

Phần này trình bày chi tiết gradient của fuzzy cross-entropy~\eqref{eq:umap-C-opt} dùng trong UMAP (Mục~\ref{sec:umap}).

Đặt $\mathbf{r}_{ij} = \mathbf{y}_i - \mathbf{y}_j$. Trước tiên tính đạo hàm của $w_{ij}$ theo $\mathbf{y}_i$:
\[
\frac{\partial w_{ij}}{\partial \mathbf{y}_i}
= \frac{\partial w_{ij}}{\partial d_{ij}} \cdot \frac{\partial d_{ij}}{\partial \mathbf{y}_i}
= \frac{-2ab\, d_{ij}^{2b-1}}{(1 + a\, d_{ij}^{2b})^2} \cdot \frac{\mathbf{r}_{ij}}{d_{ij}}
= \frac{-2ab\, d_{ij}^{2b-2}}{(1 + a\, d_{ij}^{2b})^2}\,\mathbf{r}_{ij}.
\]
Gradient của $C$ theo $\mathbf{y}_i$ trên cạnh $(i,j)$:
\[
\frac{\partial C_{ij}}{\partial \mathbf{y}_i}
= -\frac{B_{ij}}{w_{ij}}\,\frac{\partial w_{ij}}{\partial \mathbf{y}_i}
+ \frac{1-B_{ij}}{1-w_{ij}}\,\frac{\partial w_{ij}}{\partial \mathbf{y}_i}
= \underbrace{\frac{2ab\, B_{ij}\, d_{ij}^{2b-2}}{1 + a\, d_{ij}^{2b}}}_{\text{hút}}\,\mathbf{r}_{ij}
\;-\; \underbrace{\frac{2ab\,(1-B_{ij})\, d_{ij}^{2b-2}}{(1-w_{ij})(1 + a\, d_{ij}^{2b})^2}}_{\text{đẩy}}\,\mathbf{r}_{ij}.
\]
Cập nhật $\mathbf{y}_i \leftarrow \mathbf{y}_i - \eta\,\frac{\partial C_{ij}}{\partial \mathbf{y}_i}$, nên \emph{lực} (hướng cập nhật) là $\mathbf{f}_{ij} = -\frac{\partial C_{ij}}{\partial \mathbf{y}_i}$:
\begin{itemize}
  \item \textbf{Lực hút} (cặp láng giềng, $B_{ij}$ lớn)---kéo $\mathbf{y}_i$ về phía $\mathbf{y}_j$:
  \[
  \mathbf{f}_{ij}^{\mathrm{attr}} = \frac{-2ab\, B_{ij}\, d_{ij}^{2b-2}}{1 + a\, d_{ij}^{2b}}\,\mathbf{r}_{ij}.
  \]
  \item \textbf{Lực đẩy} (cặp không láng giềng, $B_{ik} \approx 0$)---đẩy $\mathbf{y}_i$ ra xa $\mathbf{y}_k$:
  \[
  \mathbf{f}_{ik}^{\mathrm{rep}} = \frac{2ab\, d_{ik}^{2b-2}}{(1-w_{ik})(1 + a\, d_{ik}^{2b})^2}\,\mathbf{r}_{ik}.
  \]
\end{itemize}

\textbf{Cập nhật SGD mỗi epoch.} Duyệt qua tất cả các cạnh $(i,j)$ trong đồ thị thưa (chỉ những cặp mà $B_{ij} > 0$). Với mỗi cạnh:
\begin{enumerate}
  \item Tung đồng xu với xác suất mặt ngửa $= B_{ij}$. Nếu ngửa: cập nhật $\mathbf{y}_i \leftarrow \mathbf{y}_i + \eta\,\mathbf{f}_{ij}^{\mathrm{attr}}$ (kéo $\mathbf{y}_i$ về phía $\mathbf{y}_j$). Cạnh có $B_{ij}$ lớn được chọn thường xuyên hơn $\to$ lực hút mạnh hơn.
  \item Chọn ngẫu nhiên \texttt{n\_neg\_samples} điểm $k$ (lấy mẫu âm---giả định các cặp ngẫu nhiên không phải láng giềng). Với mỗi $k$: cập nhật $\mathbf{y}_i \leftarrow \mathbf{y}_i + \eta\,\mathbf{f}_{ik}^{\mathrm{rep}}$ (đẩy $\mathbf{y}_i$ ra xa $\mathbf{y}_k$).
\end{enumerate}

\subsection{Khai triển gradient MLE, NCE, NEG cho nhúng láng giềng}
\label{app:nce-gradient}

Phần này khai triển tường minh ba gradient trong khung nhúng láng giềng (Mục~\ref{sec:research-contrastive}), đồng thời chứng minh $F_{\text{rep}} \propto \frac{1}{\overline{Z}}$ từ~\eqref{eq:repulsion-scaling}. Ký hiệu chung: $d_{ij} = \|\mathbf{y}_i - \mathbf{y}_j\|$, $\mathbf{r}_{ij} = \mathbf{y}_i - \mathbf{y}_j$, $\phi(ij) := \frac{1}{1+d_{ij}^2}$.

\paragraph{A. Đạo hàm kernel Cauchy.}
Bằng chain rule:
\begin{align*}
  \frac{\partial \phi(ij)}{\partial \mathbf{y}_i}
  &= \frac{d\phi}{d\,d_{ij}}\cdot\frac{\partial d_{ij}}{\partial \mathbf{y}_i}
  = \frac{-2d_{ij}}{(1+d_{ij}^2)^2}\cdot\frac{\mathbf{r}_{ij}}{d_{ij}}
  = -\frac{2}{(1+d_{ij}^2)^2}\,\mathbf{r}_{ij}
  = -2\phi(ij)^2\,\mathbf{r}_{ij}.
\end{align*}

\paragraph{B. Gradient MLE (t-SNE).}
Khai triển $\mathcal{L}_{\text{t-SNE}} = -\sum_{ij}p(ij)\log\phi(ij) + \log Z(\theta)$ (dùng $\sum_{ij}p(ij)=1$):
\begin{align*}
  \frac{\partial}{\partial\mathbf{y}_i}\bigl(-\log\phi(ij)\bigr)
  &= -\frac{1}{\phi(ij)}\cdot\bigl(-2\phi(ij)^2\,\mathbf{r}_{ij}\bigr) = 2\phi(ij)\,\mathbf{r}_{ij},\\[4pt]
  \frac{\partial \log Z(\theta)}{\partial \mathbf{y}_i}
  &= \frac{1}{Z(\theta)}\sum_{j'\neq i}\frac{\partial\phi(ij')}{\partial\mathbf{y}_i}
  = \frac{-2}{Z(\theta)}\sum_{j'\neq i}\phi(ij')^2\,\mathbf{r}_{ij'}.
\end{align*}
Kết hợp, với $q_\theta(ij) = \frac{\phi(ij)}{Z(\theta)}$:
\begin{equation}
  \frac{\partial\mathcal{L}_{\text{t-SNE}}}{\partial\mathbf{y}_i}
  = 2\sum_{j\neq i}\!\left(\frac{p(ij)}{q_\theta(ij)}-1\right)\frac{1}{Z(\theta)}
    \cdot\frac{\partial\phi(ij)}{\partial\mathbf{y}_i}.
  \label{eq:grad-tsne-app}
\end{equation}

\paragraph{C. Gradient NCE tổng quát và NC-t-SNE.}
Với $f = q_{\theta,Z}(x) = \frac{\phi_\theta(x)}{Z}$ và $c = m\xi(x)$, tính đạo hàm từng số hạng của~\eqref{eq:nce-loss}:
\[
  \frac{d}{d\theta}\log\frac{f}{f+c} = \frac{c}{f(f+c)}\frac{df}{d\theta},
  \qquad
  \frac{d}{d\theta}\log\frac{c}{f+c} = \frac{-1}{f+c}\frac{df}{d\theta}.
\]
Lấy kỳ vọng và kết hợp (với $\frac{df}{d\theta} = \frac{1}{Z}\,\frac{d\phi_\theta}{d\theta}$):
\begin{equation}
  \frac{d\mathcal{L}_{\text{NCE}}}{d\theta}
  = -\sum_x\frac{m\xi(x)}{q_{\theta,Z}(x) + m\xi(x)}\left(\frac{p(x)}{q_{\theta,Z}(x)} - 1\right)
    \frac{1}{Z}\frac{d\phi_\theta(x)}{d\theta}.
  \label{eq:nce-grad-general}
\end{equation}
Cụ thể hóa theo $\mathbf{y}_i$ cho nhúng láng giềng:
\begin{equation}
  \frac{\partial\mathcal{L}_{\text{NC-t-SNE}}}{\partial\mathbf{y}_i}
  = 2\sum_{j \neq i}\underbrace{\dfrac{m\xi(ij)}{\dfrac{\phi(ij)}{Z} + m\xi(ij)}}_{\xrightarrow{m\to\infty} 1}
    \!\left(\frac{p(ij)}{q_{\theta,Z}(ij)} - 1\right)\frac{1}{Z}
    \cdot\frac{\partial\phi(ij)}{\partial\mathbf{y}_i},
  \label{eq:grad-nctsne-app}
\end{equation}
và gradient Neg-t-SNE thu được khi thay $Z \to \overline{Z}$ (cố định):
\begin{equation}
  \frac{\partial\mathcal{L}_{\text{Neg-t-SNE}}}{\partial\mathbf{y}_i}
  = 2\sum_{j \neq i}\underbrace{\dfrac{m\xi(ij)}{\dfrac{\phi(ij)}{\overline{Z}} + m\xi(ij)}}_{\in[0.5,\,1]}
    \!\left(\dfrac{p(ij)}{\dfrac{\phi(ij)}{\overline{Z}}} - 1\right)\dfrac{1}{\overline{Z}}
    \cdot\dfrac{\partial\phi(ij)}{\partial\mathbf{y}_i}.
  \label{eq:grad-negtsne-app}
\end{equation}

\paragraph{D. Phân tách lực hút/đẩy và chứng minh $F_{\text{rep}} \propto \frac{1}{\overline{Z}}$.}
Xấp xỉ hệ số dampening bằng $1$ (hợp lệ khi $\overline{Z}$ lớn) và tách tổng~\eqref{eq:grad-negtsne-app} theo hai nhóm:
\begin{align*}
  \frac{\partial\mathcal{L}_{\text{Neg-t-SNE}}}{\partial\mathbf{y}_i}
  &\approx 2\sum_{j\neq i}\!\left(\frac{p(ij)}{\frac{\phi(ij)}{\overline{Z}}} - 1\right)\frac{1}{\overline{Z}}
    \cdot\frac{\partial\phi(ij)}{\partial\mathbf{y}_i}\\
  &= \underbrace{2\sum_{\substack{j:\,p(ij)>0}}\!\left(\frac{\overline{Z}\,p(ij)}{\phi(ij)}-1\right)
      \frac{1}{\overline{Z}}\cdot\frac{\partial\phi}{\partial\mathbf{y}_i}}_{\text{lực hút}}
  + \underbrace{2\sum_{\substack{j:\,p(ij)=0}}\!\left(-1\right)
      \frac{1}{\overline{Z}}\cdot\frac{\partial\phi}{\partial\mathbf{y}_i}}_{\text{lực đẩy}}.
\end{align*}
\textbf{Lực hút}: $\frac{\overline{Z}\,p(ij)}{\phi(ij)}\cdot\frac{1}{\overline{Z}} = \frac{p(ij)}{\phi(ij)}$---\emph{không phụ thuộc $\overline{Z}$}.

\textbf{Lực đẩy}: thay $\frac{\partial\phi}{\partial\mathbf{y}_i} = -2\phi(ij)^2\,\mathbf{r}_{ij}$:
\[
  \mathbf{f}_{ij}^{\text{rep}}
  = 2\cdot(-1)\cdot\frac{1}{\overline{Z}}\cdot\bigl(-2\phi(ij)^2\,\mathbf{r}_{ij}\bigr)
  = \frac{4\phi(ij)^2}{\overline{Z}}\,\mathbf{r}_{ij}, \qquad
  F_{\text{rep}} \propto \frac{\phi(ij)^2}{\overline{Z}} \propto \frac{1}{\overline{Z}}.
\]
Khi $\overline{Z}$ tăng (UMAP: $\mathcal{O}(n^2)$), lực đẩy suy yếu trong khi lực hút giữ nguyên---điểm tương đồng bị kéo sát nhau tạo cụm siêu co chặt, đúng hành vi quan sát được ở UMAP so với t-SNE.

\FloatBarrier
